\documentclass[11pt,a4paper]{article}
\usepackage[T1]{fontenc}
\usepackage[utf8]{inputenc}
\usepackage{lmodern,microtype,amsmath,amssymb,amsthm,mathtools}
\usepackage[margin=26mm]{geometry}
\usepackage{booktabs,longtable,array,enumitem,xcolor,fancyhdr}
\usepackage[hyphens]{url}
\usepackage[colorlinks=true,linkcolor=blue!45!black,citecolor=blue!45!black,urlcolor=blue!45!black]{hyperref}
\usepackage{listings}
\hypersetup{pdftitle={Shared coefficient positions and integral Pell carriers},
pdfauthor={Research contribution prepared for Vladimir Reshetnikov},
pdfsubject={Explicit Diophantine equations for Con(ZFC)}}
\lstset{basicstyle=\ttfamily\small,breaklines=true,columns=fullflexible,keepspaces=true}
\setlength{\emergencystretch}{3em}
\setlength{\parskip}{4pt}
\setlength{\parindent}{0pt}
\setlist{nosep,leftmargin=*}
\newcommand{\Z}{\mathbb Z}
\newcommand{\N}{\mathbb N}
\newcommand{\Zp}{\mathbb Z_{>0}}
\newcommand{\Con}{\operatorname{Con}}
\newcommand{\ZFC}{\mathsf{ZFC}}
\newcommand{\degp}{\operatorname{deg}}
\newtheorem{theorem}{Theorem}[section]
\newtheorem{lemma}[theorem]{Lemma}
\newtheorem{proposition}[theorem]{Proposition}
\newtheorem{corollary}[theorem]{Corollary}
\theoremstyle{definition}\newtheorem{definition}[theorem]{Definition}
\theoremstyle{remark}\newtheorem{remark}[theorem]{Remark}
\pagestyle{fancy}\fancyhf{}
\fancyhead[L]{\small Smaller Diophantine consistency equations}
\fancyhead[R]{\small 23 October 2026}
\fancyfoot[C]{\thepage}
\setlength{\headheight}{14pt}
\title{Shared coefficient positions and integral Pell carriers\\[4pt]
\large Two reductions for explicit Diophantine equations expressing $\Con(\ZFC)$}
\author{Research contribution prepared for Vladimir Reshetnikov\\
\small AI-assisted mathematical research for the ProveIt repository}
\date{23 October 2026}
\begin{document}
\maketitle
\begin{abstract}
We investigate simpler fixed integer polynomials whose absence of zeros
in their stated integer domains is provably equivalent in ZFC to syntactic
consistency of ZFC.
The comparison uses ProveIt commit
\texttt{beb0192cb8dc36f78c369c080f2adfbc1887a486}, including the six
wave-eight packages that are newer than the collection overview.
Two complementary changes are developed. Compatible equal coefficient
contributions may share an encoding position, reducing the degree and
arithmetic cost of the nine-positive-variable construction, and also improving
the ten-unrestricted-integer construction. A second reduction
uses one integral Pell root to recover, in order, an index congruence and
a missing scale divisibility; it gives a degree-26 equation in fourteen
positive variables. All coefficients are fixed by explicit finite recipes.
The archive contains literal arithmetic circuits, coefficient data,
proofs, and independently written finite checks. The new reductions are
ordinary mathematical proofs using named inherited compiler results;
they do not constitute a new end-to-end proof-assistant certificate or a
proof of global minimality.
\end{abstract}
\begin{center}
\small
\textbf{Principal constructions}\\[6pt]
\begin{tabular}{@{}lrrr@{}}\toprule
Witness domain & Unknowns & Exact degree & Binary operations\\\midrule
Strictly positive integers & 9 & 1,675,396 & 458\\
Unrestricted integers & 10 & 7,428,120 & 525\\
Strictly positive integers & 14 & 26 & 92\\\bottomrule
\end{tabular}
\end{center}
\clearpage
\tableofcontents
\clearpage
\section{Results and the meaning of simplicity}
\label{sec:results}
The intended logical statement is
\begin{equation}\label{eq:target}
 \ZFC\vdash
 \bigl(\neg\exists\mathbf x\in\Zp^n\;P(\mathbf x)=0\bigr)
 \;\longleftrightarrow\;\Con(\ZFC).
\end{equation}
Here $\Con(\ZFC)$ means that there is no finite derivation of falsity from
the fixed syntactic ZFC axiom predicate used by ProveIt. No consistency
assumption enters the construction. In particular, proving the equivalence
does not prove either side of it.

There is no agreed total order on polynomial complexity. We report the
witness domain, number of witnesses, exact total degree, arithmetic circuit
size, and coefficient binding separately. A circuit operation is one live
binary addition, subtraction, or multiplication, including multiplication
by a fixed numeral. Reusing a register costs nothing. Each fixed power is
implemented by ordinary multiplication gates. Coefficient leaves are fixed
integers, and their bit lengths are a separate cost.

\begin{table}[htbp]\centering\small
\begin{tabular}{@{}lrrr@{}}\toprule
Construction & Positive variables & Exact degree & Operations\\\midrule
Package 59, recommended nine-variable form~\cite{p59} & 9 & 1,910,404 & 473\\
Shared positions, primary form & 9 & 1,675,396 & 458\\
Shared positions, cheaper tradeoff & 9 & 1,738,228 & 457\\
Package 59, fourteen-variable form~\cite{p59} & 14 & 52 & 79\\
Integral carriers, primary form & 14 & 26 & 92\\
Integral carriers, cheaper tradeoff & 14 & 30 & 85\\
Package 58, degree-24 fifteen-variable form~\cite{p58} & 15 & 24 & 90\\
Package 51, 66-operation positive circuit~\cite{p51} & 15 & 157 & 66\\
\bottomrule\end{tabular}
\caption{Comparisons at the pinned repository snapshot. Every row is a
single polynomial equation over strictly positive integers. Earlier rows
are cited constructions, not newly rerun end-to-end semantic proofs.}
\label{tab:main}
\end{table}

The primary nine-variable change lowers degree by approximately $12.30\%$
and operation count by approximately $3.17\%$ relative to package 59.
The fourteen-variable degree falls by $50\%$, at the cost of thirteen
additional gates. The degree-24 fifteen-variable equation and the
66-operation equation retain their distinct advantages. Neither new
family reduces the smallest variable count in the inspected collection.
These are concrete improvements to that collection, not a claim that the
longstanding problem of the globally simplest consistency equation has
been solved.

Over ten unrestricted integer variables, the shared-position adapter has
exact degree $7{,}428{,}120$ and $525$ operations. It improves on package~59's
$8{,}469{,}528$ at $533$ operations, and on the degree of package~57's
$8{,}317{,}400$ at $544$ operations~\cite{p57,p59}. The older package~54
point uses fewer operations ($513$), at degree $25{,}838{,}224$~\cite{p54}; we therefore
do not claim the least arithmetic cost in that domain. This third
construction uses the integer-domain finalizer based on Sun's
ten-unknown theorem~\cite{sun}, with its hypotheses and bootstrap inherited
from the pinned package~59. In its version of~\eqref{eq:target}, the
quantified domain is $\Z^{10}$.

The constructions deliberately retain two different complete coefficient
bindings. The nine-variable and ten-integer families use the original source-31 tuple;
the fourteen-variable family uses the compact source-44 tuple. Coefficient
ports with the same names must not be exchanged between them. Their
bindings, hashes, and regeneration instructions accompany the circuits.

\section{Logical scope and inherited interfaces}
\label{sec:scope}
The existing ProveIt pipeline relates a literal finite computation to a
syntactic ZFC refutation, translates that computation into arithmetic,
and encodes the arithmetic using Pell equations or coefficient extraction.
Our work replaces specific arithmetic reductions within that pipeline.
The proofs below establish existence equivalences with the supplied source
systems. Composing those equivalences with the cited native compiler and
refutation calibration gives~\eqref{eq:target}.

Three levels of evidence must be kept distinct. The mathematical proofs
establish universal implications over integers. The finite programs check
literal circuit transcription, all finite source identities, exact-degree
certificates, coefficient hashes, and regeneration. A proof-assistant
kernel certificate for a complete closed consistency equation is a
separate artifact. ProveIt's active research records such closed
constructions, but those theorems do not automatically certify the new
literal polynomials here. We do not supply a newly checked object-language
ZFC derivation or a complete new Lean/Rocq build.

All operations needed by the new proofs are elementary arithmetic,
finite polynomial identities, the stated Pell results, and the explicitly
identified inherited computation theorem. They are formalizable in ZFC.
The imported files retain their own statements of scope and provenance.
The final review records any remaining reliance on those interfaces.

\section{Why shared coefficient positions are valid}
\label{sec:sharing_principle}
A useful way to encode several equations is to extract chosen coefficients
from a product of polynomials. Let
\[
 Z(B,\mathbf x)=\sum_{i=1}^{m}x_iB^{w_i},\qquad
 A_d(B)=\sum_{u}a_{d,u}B^u.
\]
At a chosen target $t_j$, the degree-$d$ contribution is
\begin{equation}\label{eq:coefficient_identity}
 [B^{t_j}]A_d(B)Z(B,\mathbf x)^d
 =\sum_{\substack{u+\mathbf w\cdot\boldsymbol\alpha=t_j\\
             |\boldsymbol\alpha|=d}}
   a_{d,u}\binom d{\boldsymbol\alpha}\mathbf x^{\boldsymbol\alpha}.
\end{equation}
The coefficient ring can itself be a polynomial ring in named fixed
coefficient ports. Thus equality of two coefficient contributions can be
checked symbolically before enormous numerals are substituted.

\begin{proposition}[Sharing a compatible coefficient position]
Suppose the right side of~\eqref{eq:coefficient_identity} equals the required
source polynomial for each target and each homogeneous degree. Then the
coefficient decoder is correct, regardless of whether a kernel position
serves more than one target. When several requirements at the same position
ask for the identical coefficient, that coefficient is stored once.
\end{proposition}
\begin{proof}
Expand the finite product and collect powers of $B$. Formula
\eqref{eq:coefficient_identity} enumerates every contribution, including
all interactions with every monomial in $Z^d$. A kernel position contains
one coefficient. Each target sees that coefficient whenever the exponent
condition holds. Therefore the checked identities establish the decoder
directly; no assumption that every target has private kernel positions is
needed. Storing repeated requirements by addition would instead change
the coefficient and is not the construction asserted here.
\end{proof}

This is a relaxation of the former placement constraint. It does not
permit arbitrary collisions. The independent checker enumerates every
relevant interaction and compares complete source polynomials. The
following section supplies the finite certificate, the complete polynomial
definition, and the transport of the source bounds to the final equation.

\section{Nine variables: compatible kernel coefficients}
\label{sec:shared-taps}

This contribution changes the source encoding in package~59 while
retaining its arithmetic finalizer. It gives two explicit integer
polynomials on nine strictly positive integer inputs:
\begin{center}
\begin{tabular}{lrrr}
\toprule
Construction & Exact degree & Binary operations & Source height\\
\midrule
Package 59, inherited comparator & $1{,}910{,}404$ &473&2161\\
Shared coefficients, primary & $1{,}675{,}396$ &458&1873\\
Shared coefficients, economy & $1{,}738{,}228$ &457&1950\\
\bottomrule
\end{tabular}
\end{center}
The primary lowers degree by approximately $12.30\%$ and operation count
by $3.17\%$ against this comparator. The two new points are mutually
nondominated. They retain the existing variable count, the entire original
native coefficient tuple, and the same unbounded compiler and ratio-theorem
dependencies. These results constitute a measured local improvement; they
do not establish a global minimum.

\subsection{The shared-coefficient lemma}

Let $r_1,\ldots,r_{22}$ be the original quadratic source residuals in
$31$ positive coordinates, specified by the literal $65$-gate syntax in
\texttt{nine/data/source\_system65.json}. Put
\[
 \rho_{14}=r_{14}+r_2,\qquad \rho_{15}=r_{15}+r_3,
 \qquad \rho_j=r_j\quad(j\ne14,15).
\]
These row operations are invertible over the integers. For a monomial
$p^\alpha$ of ordinary degree $d\le2$ with coefficient $a_{j,\alpha}$ in
$\rho_j$, define
\[
 c_{j,\alpha}=\begin{cases}
 2a_{j,\alpha},&d=0,1,\\
 a_{j,\alpha}\prod_i\alpha_i!,&d=2.
 \end{cases}
\]
Coefficients of absent monomials are zero. Choose positive distinct
weights $w_i$ and distinct targets $q_j$.

\begin{lemma}[Compatible kernel coefficients]
Suppose that, for $d=0,1,2$, a finitely supported table $k_d$ on
nonnegative integers satisfies
\[
 k_d(q_j-\alpha\cdot w)=c_{j,\alpha}
 \tag{\(*\)}
\]
for every degree-$d$ monomial, with $k_d(e)=0$ for $e<0$. Define
\[
 D_d(T)=\sum_e k_d(e)T^e,\qquad U_p(T)=\sum_i p_iT^{w_i}.
\]
Then
\[
 [T^{q_j}]\bigl(D_2(T)U_p(T)^2+D_1(T)U_p(T)+D_0(T)\bigr)
   =2\rho_j(p).
\]
Repeated requirements in \((*)\) may share a position whenever their
complete coefficient polynomials agree. The shared coefficient occurs
once in $D_d$.
\end{lemma}
\begin{proof}
The coefficient of $p^\alpha$ in the target is
$k_d(q_j-\alpha\cdot w)$ times the multinomial coefficient in $U_p^d$.
For $d=2$ the latter is $1$ on a square and $2$ on a cross product, giving
$2a_{j,\alpha}$ with the stated weighting. Degrees zero and one are
immediate. Checking absent monomials excludes every unwanted contribution.
\end{proof}

The earlier construction prohibited every same-degree overlap of source
records. The lemma replaces this prohibition by exact compatibility. It
is essential to store an agreed coefficient once: summing its repeated
requirements changes the polynomial and invalidates the target identity.

\subsection{Explicit finite data and source bounds}

Both certificates use the unchanged weight vector, in source-variable order,
\[
\begin{split}
w={}&(1,7,61,36,16,85,20,67,43,11,34,13,96,80,29,50,\\
 &2,109,110,112,116,91,53,55,27,101,92,9,107,37,60).
\end{split}
\]
The primary targets in row order $1,\ldots,22$ are
\[
\begin{split}
q={}&(1117,28,1442,1045,1841,1330,1819,614,267,1384,1379,\\
 &889,337,1592,413,673,1115,772,1107,1413,1031,85).
\end{split}
\]
Its $89$ weighted requirements use $84$ distinct kernel coefficients. The
nontrivial groups are as follows; coordinate names are the original source
names, not the nine final inputs.
\begin{center}
\begin{tabular}{rrlp{5.4cm}}
\toprule
Degree&Position&Value&Shared requirements\\
\midrule
1&1324&$2$&Row 11: $C2$; row 10: $\mathit{Delta}$\\
2&39&$-1$&Row 22: $g\,\mathit{Baux}$; row 9: $ca$\\
2&588&$2$&Row 8: $\tau^2$; row 18: $V^2$\\
2&1032&$-1$&Row 4: $sN$; row 19: $i\,C2$; row 17: $fo$\\
\bottomrule
\end{tabular}
\end{center}
The complete kernel tables are in \path{nine/data/kernels.json}.
The economy tables are separate files.
The checker derives the residuals independently from the original source
syntax and verifies all
$22(1+31+496)=11{,}616$ coefficient requirements for each certificate,
including absent monomials. Thus this finite certificate verifies universal
polynomial identities rather than agreement at sampled assignments.

Write $W=\max w_i$, $Q=\max q_j$, and
$H=\max_d(\deg D_d+dW)$. The exact values are
\begin{center}
\begin{tabular}{lrrrrr}
\toprule
Variant&$W$&$Q$&$H$&$(\deg D_0,\deg D_1,\deg D_2)$&Unique taps\\
\midrule
Primary&116&1841&1873&$(1384,1739,1641)$&84\\
Economy&116&1918&1950&$(1508,1816,1718)$&86\\
\bottomrule
\end{tabular}
\end{center}
Both quadratic kernels have leading coefficient $+1$.

Use the complete original31 tuple, whose port bit bounds are
\[
 (1064921,1070004,1064921,1064922).
\]
Its compressed payloads, builder,
transition input and inherited notices accompany the article. Every
weighted kernel coefficient has absolute value below
$M_c=2^{2134927}$. Define
\[
 \kappa=2134939,\qquad a=2^\kappa=4096M_c,\qquad B=ab^2.
\]
For positive $b,g$, set
\begin{align*}
 \ell&=\sum_i B^{w_i}, & U_s&=\ell+g,
 &G_s&=(D_2(B)U_s+D_1(B))U_s+D_0(B),\\
 T_*&=\sum_jB^{q_j}, & N_0&=B^{117},
 &Z&=B^{H+1},\\
 M_*&=N_0-1-(b-1)\ell,
 &S_*&=g+N_0(2G_s+T_*),
 &T_c&=M_*+N_0(B-2)T_*.
\end{align*}
Also put $t=bB^{116}$ and $G_{\rm cap}=32t$.

Here $U_s\ge B^{116}$, and
$D_2(B)>\tfrac12B^{\deg D_2}$. Relative to
$B^{\deg D_2}U_s^2$, the possibly negative terms have absolute values
less than $2M_c/B^{18}$ and, respectively, $2M_c/B^{489}$ or
$2M_c/B^{442}$. Both are less than $1/8$. Hence $G_s>0$ before digit
recovery. On the ratio-derived range $g<4bB^{116}$, one has $U_s<6bB^{116}$
and
\[
 \frac{2G_s}{Z}<\frac{86M_cb^2}{B-1}<\frac1{32}.
\]
Since $H-Q=32$, both $T_*$ and $(B-2)T_*$ are smaller than $Z/64$.
Together with $g<N_0$, this gives $S_*+T_c<B^{H+118}$.

When $b$ is a power of two and $S_*\mathbin{\&}T_c=0$, the low mask recovers
$g=\sum_i z_iB^{w_i}$ with $0\le z_i<b$, and $U_s$ encodes
$p_i=z_i+1>0$. Every convolution coefficient has absolute value below
\[
 M_c\bigl((\textstyle\sum_i p_i)^2+\sum_i p_i+1\bigr)
 \le 993M_cb^2<B/4.
\]
The inherited sparse-borrow lemma therefore turns $S_*\mathbin{\&}T_c=0$
into exactly the common source zero condition. Compatible sharing does not
alter this bound: it concerns unique kernel coefficients. Conversely, every
source zero works for sufficiently large admissible bases; the all-one
source tuple fails row 5, so $g>0$, and $g<bB^{116}$ gives $32g<G_{\rm cap}$.

\subsection{Binary calibration and the complete final equation}

Let $E_c=H+118$ and $A_c=H+255$. Define the fixed integers
\begin{align*}
 \delta&=(A_c-15)\bmod2,&
 \lambda&=(22-A_c\kappa)\bmod2,\\
 n_B&=(A_c-15-\delta)/2,&
 n_2&=((15+\delta)\kappa+\lambda-22)/2.
\end{align*}
Then
\[
 L=2^\lambda bB^{E_c},\qquad
 R=L(S_*+T_c+1)-S_*-1,\qquad N=2^{n_2}b^\delta B^{n_B}.
\]
The primary parameters $(E_c,\lambda,\delta,n_B,n_2)$ are
$(1991,0,1,1057,17079501)$; the economy parameters are
$(2068,1,0,1096,16012032)$.
If $b=2^k$, $B=2^v$, $L=2^u$, the two-block population identity is
\[
 \operatorname{pop}(R)=u+\operatorname{pop}(T_c)-\tau,
 \qquad \operatorname{pop}(T_c)=139v-31k-22,
\]
where $\tau\ge0$ vanishes exactly for disjoint source words. The parameters
give $u+\operatorname{pop}(T_c)=2\log_2N$. Legendre's valuation formula
therefore gives $N^2\mid\binom{2R}{R}$ exactly when the source words are
disjoint. Both parity corrections are part of the definition.

The final inputs are $(b,g,h,s,w,C,i,j,n)\in\mathbb Z_{>0}^9$;
put $j_0=j-1$. The complete remaining registers are
\begin{align*}
 U&=N^2w,&Y&=N^2s,\\
 M_0&=RY,&P&=2M_0^2U,\\
 K&=R+1+h(P-1),& A&=M_0(U+1),\\
 B_1&=2R+1,& W_1&=bw,\\
 J&=4A-5,&Q_1&=3W_1C-2(W_1^2-1),\\
 D&=(A^2-1)C^2+1,& k_0&=iJC,\\
 E_0'&=k_0DC,&F&=(A^2-1)(E_0')^2+1,\\
 G&=A+F(CD+1-A),&H_0&=C+F(B_1+j_0C),\\
 I&=(G^2-1)H_0^2+1,\\
 a_0&=(P^2-1)K^2+1,&a_1&=DFI,\\
\end{align*}
\[
 m_0=K^2(4t-g)-16t(C-KY)^2.
\]
\begin{align*}
 V_0&=PK,&V_1&=k_0(8D^2-4D-1)GH_0,\\
 \Omega&=Jm_0,& E_0&=J(n+W_1^2)+Q_1,\\
 \Theta&=E_0-\Omega(V_0+V_1).
\end{align*}
The desired single polynomial is
\begin{equation}\label{eq:new-nine}
 \boxed{\mathcal F_9=
  [\Theta^2-\Omega^2(a_0+64a_1)]^2
      -256\Omega^4a_0a_1.}
\end{equation}
Every power has a fixed integer exponent, and every register is an
abbreviation. The literal evaluator and every nonliteral fixed leaf are
supplied. Each variant uses its own twelve-leaf binding; the native tuple is
shared whole.

\begin{proposition}
For each of the two fixed parameter sets, equation~\eqref{eq:new-nine}
has a positive integer zero exactly when the original fixed native source
has a positive common zero. Composing with the inherited source compiler
and calibration gives the stated ZFC-provable consistency equivalence.
\end{proposition}
\begin{proof}
The compatible-coefficient lemma, the positivity and carry bounds, and the
exact population calculation prove the same source contract as package~59.
The arithmetic finalizer above is unchanged. Its proof, retained in
\texttt{nine/inherited/nine\_free\_proof.md}, uses only that contract and
the numerical conditions $R,N\ge8$, $b\le R,N$, which hold here.
Its strict majorants first force $m_0>0$ from a norm zero. That recovers the
carrier range before radical separation is invoked. Both radicals are then
integral, $J\mid Q_1$ follows, and the Jones ratio theorem yields the binary
base and binomial divisibility. The source contract recovers a common zero.
Conversely, a source zero admits arbitrarily large admissible bases; the
ratio theorem and exact CRT substitution supply the auxiliary witnesses,
and its positive quotient construction supplies $n>0$. No step uses
separate ownership of a shared kernel coefficient or the old height 2161.
\end{proof}

\subsection{Degree cancellation and reproducible evidence}

Raw maximum-degree propagation through the norm is insufficient. Put
$r=\deg R=4E_c-1$ and $q=\deg N=2n_B+\delta$.
With $T_p=k_0(8D^2-4D-1)$, use the exact identities
\begin{align*}
 T_p^2-64DF&=(8D+1)k_0^2-64D,\\
 \Delta_1:=V_1^2-64a_1
  &=\bigl((8D+1)k_0^2-64D\bigr)(GH_0)^2+64DF(H_0^2-1),\\
 \Theta^2-\Omega^2(a_0+64a_1)
  &=E_0^2-2E_0\Omega(V_0+V_1)
     +\Omega^2\bigl(K^2-1+\Delta_1+2V_0V_1\bigr).
\end{align*}
The unique leading term in the last line is $\Omega^2\Delta_1$, of degree
$52r+200q+4W+158$. The final subtracted term in
\eqref{eq:new-nine} has degree $76r+288q+8W+216$, strictly below twice that
quantity. Therefore
\[
 \deg\mathcal F_9=104r+400q+8W+316.
\]
For $(r,q)=(7963,2115)$ this is $1{,}675{,}396$.
The diagonal leading coefficient is
$2^{186}a^{837288}c^{400}>0$, where $c=2^{17079501}$, proving that the
upper bound is attained with the actual fixed tuple. For the economy
parameters $(r,q)=(8271,2192)$ the degree is $1{,}738{,}228$.

The independent checker verifies the exact source identities, literal
regeneration, actual coefficient payloads, operation recounts, the resolved
leading forms and the strict degree comparison. It also compares the DAG
with an independently written direct formula at modular assignments;
these are transcription checks. The larger degree-certificate DAGs are not
the evaluators counted as 458 or 457 gates. The translations to nonnegative
coordinates require seven additional gates, giving 465 and 464.

All checks run with either ordinary Python or \texttt{python3 -O}; altered
shared coefficients fail the mathematical identity checks. These checks
are distinct from an end-to-end proof-assistant certificate and from an
object-language ZFC derivation. The new finite source lemma is a natural
next formalization target, followed by composition with the existing
compiler endpoint. Joint optimization of row operations, coefficient
equality classes and target positions is a separate further research
question.


\section{Two divisibilities and an index quotient in one Pell root}
\label{sec:pell-carriers}

This section gives a second, direct route to smaller equations. It uses the
authentic compact native tuple and ordinary Pell kernel, and does not use
Jones's nine-variable compression. The new result has fourteen strictly
positive unknowns, exact degree $26$, and a $92$-operation evaluator. A
companion has fourteen positive unknowns, degree $30$, and $85$ operations.
The inspected wave-eight fourteen-positive point is degree $52$ at $79$
operations. Thus the first construction halves its degree and adds thirteen
gates; the second uses six additional gates and reduces its degree by
$42.3\%$. These are tradeoffs within the pinned collection, with no claim
of a global simplicity minimum.

The structural step combines the integral-root divisibility carrier of
package~58 and the main-root index absorption of package~59. A single
existing Pell root now enforces a weak congruence first, determines the
Pell index, and only then enforces a stronger congruence restoring another
divisibility. This ordered recovery permits the coordinate $X$ to be
supplied directly while deleting the old index quotient $h$.

\subsection{Fixed coefficients and exact equation}

Use the \emph{whole} package-44 compact tuple, whose binding is included
under \texttt{pell/coefficients/}. Its identifier is
\texttt{wang943-complement67-shared-anchor-length8807-prime581197}.
Set
\[
 B=2^{590069},\qquad \beta=B-1,\qquad
 K_0=D_C+BD_R,\qquad M_F=M_{F,\mathrm{native}}+B-1.
\]
The fourth port is $M_C$. The high mask $M_F$ is already shifted. The four
little-endian binary integers occupy $295{,}388$ raw bytes and $9{,}445$
gzip bytes. Their recipes, payloads, exact hashes, literal machine data,
native adapter, and independent coefficient auditor accompany the paper.
All four ports are fixed numerals, not quantified coordinates.

The fourteen positive inputs are
\[
 j,F,\alpha,z,f,i,T,x,Y,r,\eta,\zeta,u,\gamma.
\]
Define polynomial abbreviations
\begin{align*}
 J&=3j+1,& q&=\beta J+1,& C&=q-F-\alpha,&X&=x+q,\\
 E&=XY,& k&=\eta+\zeta,&c&=kY+\eta,& a&=E+Y,\\
 H&=4a+3,&\Delta&=a^2+H=(a+2)^2-1,&\ell&=EkY.
\end{align*}
Also set
\begin{align*}
 R&=\bigl[(q-1)(q-F)+\alpha\bigr](q^2-1)
                      +(M_C+qM_F)J,\\
 S&=4qY,\qquad M=qS=4q^2Y,\qquad v=k-R-1,\\
 D_*&=M(X+ac)+H(M\gamma-v+SX),\\
 U&=\Delta(f^2-1),\qquad I=ic^2,\qquad V=f^2(cT-R)-c.
\end{align*}
The five residuals are
\begin{align}
 A_0&=[q^3(\ell+r)+Y]^2-q^6[\ell(\ell+k)+1],
                                                    \label{eq:pell-first}\\
 A_m&=D_*^2-M^2(\Delta c^2+1),                       \label{eq:pell-main}\\
 A_t&=q[K_0C-F-z(q-1)]+XC,                          \label{eq:pell-transport}\\
 A_a&=(V+u)^2-Uu(2V+u)-1,                           \label{eq:pell-aux}\\
 A_s&=I^2-U.                                       \label{eq:pell-strong}
\end{align}
The single equation is
\begin{equation}
 \boxed{P_{14,26}=A_0^2+A_m^2+A_t^2+A_a^2+A_s^2=0.}
 \label{eq:pell26}
\end{equation}
There are no variable exponents, radicals, divisions or extra quantified
registers in this definition. The supplied literal circuit is
\texttt{pell/circuits/p14\_degree26.json}.

\begin{theorem}[Fourteen-positive equation]
\label{thm:pell26}
For the stated authentic coefficient tuple, $P_{14,26}$ has exact total
degree $26$ and a $92$-gate evaluator. The new reductions below, composed
with the explicitly supplied ordinary native compiler and calibration,
establish by ordinary mathematics formalizable in ZFC that
\[
 \ZFC\vdash
 \bigl[\neg\exists\mathbf t\in\Zp^{14}\;
                   P_{14,26}(\mathbf t)=0\bigr]
       \longleftrightarrow\Con(\ZFC).
\]
This is not a newly kernel-checked end-to-end formalization or a supplied
derivation file in the object-language ZFC calculus.
\end{theorem}

\subsection{The two-stage carrier lemma}

The elementary fact behind the construction is that a rational square root
of an integer is integral: $m^2\mid n^2$, with $m>0$, implies $m\mid n$.
At a zero of $A_0$, it gives
\[
 q^3\mid Y,\qquad s=Y/q^3>0,\qquad
 \tau=\ell+r+s>0,\qquad \tau^2=\ell(\ell+k)+1.
\]
Since $590069$ is odd, $\beta\equiv1\pmod3$. Thus the supplied
$J=3j+1$ gives $q\equiv2\pmod3$. As $q\mid Y$,
\[
 H=4Y(X+1)+3\equiv3\pmod q,
 \qquad \gcd(q,H)=1.
\]
The odd integer $H$ also satisfies $\gcd(Y,H)\mid3$. Hence
\begin{equation}
 \gcd(S,H)\mid3,\qquad
 \mu:=\frac S{\gcd(S,H)}\ge\frac{4qY}{3}.
 \label{eq:pell-mu}
\end{equation}
For $Y\ge q^3$ and $R<q^4-2$, this yields $\mu>R+1$.
Moreover both $S$ and $M$ divide $4XY^2$; for $M$ the quotient is
$XY/q^2=Xsq$, which is already integral without assuming $q\mid X$.

At a zero of $A_m$, the same rational-root lemma gives the integer
$g=D_*/M$. Reducing $Mg=D_*$ only modulo $S$ gives
\begin{equation}
 S\mid H(k-R-1),\qquad \mu\mid k-R-1.
 \label{eq:pell-weak}
\end{equation}
The Pell argument below uses this weaker congruence to recover the exact
index and show $M\mid v$. Only then do we return to the full modulus:
\begin{equation}
 M\mid HSX\quad\Longrightarrow\quad q\mid HX
       \quad\Longrightarrow\quad q\mid X.
 \label{eq:pell-strongdiv}
\end{equation}
The order of these steps prevents a circular use of the quotient $X/q$.

\subsection{Soundness: recovering the indices before the quotients}

At a zero of \eqref{eq:pell26}, all five residuals vanish. Restore $s$ and
$\tau$ as above. The cleared transport equation becomes the rational
equality
\[
 (K_0+X/q)C=F+z(q-1).
\]
Both its right side and its coefficient of $C$ are positive. Therefore
$C>0$ before $X/q$ has been shown integral. The authentic mask bounds give
\begin{gather}
 q\ge16,\qquad q(q^2-1)<R<q^4-q^3+q^2-1<q^4-2,\label{eq:pell-bounds}\\
 R-2>3q+1,\qquad E>q^4>R+2,\qquad a>R+2,
 \qquad a+2<2E.\nonumber
\end{gather}
For clarity, the bounds on $R$ follow from
\[
 G=(q-1)(q-F)+\alpha=q\alpha+(q-1)C,
 \quad q\le G\le q^2-q,
\]
and $0<(M_C+qM_F)J<(q-1)(1+2q)$. The supplied shift $X=x+q>q$
is what proves $E>q^4$ at this stage. The inequalities
$kY<c<k(Y+1)$ follow directly from $\eta,\zeta>0$.

The integer $g=D_*/M$ is positive. Indeed, $v<k$,
$b_0=X+ac>akY$ and $H<5a$ give
\[
 D_*>Mb_0-Hk>ak(MY-5)>0.
\]
Hence $g^2-\Delta c^2=1$. The first and main norms give positive Pell
indices $n,p$ with
\begin{align*}
 \tau&=\chi_{P}(n),&k&=2\psi_{P}(n),&P&=2XY^2+1,\\
 g&=\chi_A(p),&c&=\psi_A(p),&A&=a+2.
\end{align*}
Here $\chi_A(t)+\psi_A(t)\sqrt{A^2-1}
=(A+\sqrt{A^2-1})^t$. The first norm has fundamental integer-coordinate
solution $(2XY^2+1,2)$ for discriminant $XY^2(XY^2+1)$: a second
coordinate $1$ is excluded by consecutive squares. Since $P>A$ and $c>k$,
monotonicity implies $p>n$.

The Pell recurrence gives $k\equiv2n\pmod{4XY^2}$. Since $S\mid4XY^2$,
\eqref{eq:pell-weak} implies
\[
 2n=R+1+t\mu\quad(t\in\mathbb Z).
\]
The inequalities $0<R+1<\mu$ and $n>0$ force $t\ge0$. Thus $n>24$ and
$p\ge26$. Consequently
\begin{equation}
 c>A\Delta^2,\qquad c>2p,\qquad c>2R.
 \label{eq:pell-rankbounds}
\end{equation}
For the first estimate it suffices that $p\ge6$ and
$\psi_A(6)=32A^5-32A^3+6A>A(A^2-1)^2$. The second is elementary Pell
growth, and the third follows from $c\ge\psi_A(2)=2A>2R$.

The equation $A_s=0$ gives $U=I^2$ with $I=ic^2>1$, hence $f\ge2$.
The \emph{ordinary strong-rank lemma} supplied in
\texttt{pell/upstream/ORDINARY\_KERNEL.md}, Sections 4--5, now applies
with all its hypotheses established. It gives an index $m$ such that
\begin{equation}
 f=\chi_A(m),\quad p\mid m,\quad c\mid m,\quad m>2p,
 \quad I=\Delta\psi_A(m),\quad c^2\mid f^2-1.
 \label{eq:pell-strongrank}
\end{equation}
The last divisibility is a conclusion of the rank lemma; it cannot be
obtained by informally cancelling $\Delta$ in $I^2=\Delta(f^2-1)$.
The copied proof handles the possibly nonfundamental squarefree Pell
parameter and is an explicit inherited dependency.

The auxiliary equation is exactly
\[
 I^2V^2-(I^2-1)(V+u)^2=1.
\]
It excludes $V=0$ and $V+u=0$. Put $y=|V+u|>0$. The supplied quotient
and \eqref{eq:pell-strongrank} imply
\[
 V\equiv-c\pmod f,\qquad V\equiv-R\pmod c.
\]
Also $c^2\mid I$. The signed auxiliary-index lemma in the same ordinary
source gives an auxiliary index congruent to $\pm p$ and $\pm R$ modulo
$c$. The bounds $0<p,R<c/2$ force $p=R$. This argument does not assume
$V>0$. Its ordinary source quotient is restored by $T_{\rm old}=fT>0$.

Now $n<p=R$, so $|2n-R-1|<R<\mu$. Together with the congruence modulo
$\mu$ this yields $2n=R+1$. Therefore
\[
 v=k-R-1=k-2n
\]
is a positive multiple of $4XY^2$. Positivity follows from $n\ge2$ and
$\psi_P(n)>n$. In particular $M\mid v$, and
\eqref{eq:pell-strongdiv} proves $q\mid X$. Restore
\begin{equation}
 w=X/q>0,\qquad h=v/E>0,\qquad
 \gamma_{\rm old}=\gamma-v/M+w.
 \label{eq:pell-restore}
\end{equation}
These are integers. The main root becomes
$g=X+ac+H\gamma_{\rm old}$. Since $n\ge2$ gives $k\ge4P>X$, one has
$c>X$; and $g>\sqrt\Delta c>(a+1)c$. Thus
$\gamma_{\rm old}>(c-X)/H>0$. The transport equation is now the genuine
integer equality $(K_0+w)C=F+z(q-1)$, corresponding to old source
coordinate $z_{\rm old}=z+1$.

All positive ordinary source coordinates are recovered and all six source
factors equal $1$. The supplied native compiler and calibrated machine
endpoint therefore yield a finite refutation of the specified ZFC calculus.

\subsection{Completeness and the positive replacement witnesses}

Conversely, a finite refutation gives canonical normalized witnesses from
the authentic compact adapter. Transfer the strong coordinate by
$i=\Delta i_{\rm normalized}$, as in package~58. Then
$I^2=\Delta(f^2-1)$, $c\mid I$, and $f^2\equiv1\pmod c$.
Canonical indices satisfy
\[
 p=R=2n-1,\qquad n\ge3,\qquad R\equiv3\pmod4,
 \qquad X=2^R,\qquad Y=sq^3,
\]
and $k>XY$. These are inherited compiler/Pell facts, not conclusions of a
finite numerical search.

The adapter supplies $q=B^N$ and $J=1+B+\cdots+B^{N-1}$ with odd
$N=581197^e>1$. Since $B\equiv-1\pmod3$, $J\equiv1\pmod3$ and
$j=(J-1)/3>0$. Canonical $X=2^R$ with $R>3q$ gives $x=X-q>0$.
The source proves $w>q>F$, which makes $z_{\rm old}\ge2$; hence the
transport shift $z=z_{\rm old}-1$ is positive.

Supply $Y$ directly and set $r=\tau-\ell-s$. The first norm and
$\ell>k$ give $\tau<\ell+k$ and
\[
 \tau-\ell=\frac{\ell k+1}{\tau+\ell}
       >\frac{k}{3}>\frac{XY}{3}>s.
\]
Thus $r>0$ and $A_0=0$. The positive multiple $v=k-2n$ of $4XY^2$
satisfies $M\mid v$. Set
\[
 \gamma=\gamma_{\rm old}+v/M-w.
\]
This is positive for a particularly simple reason:
\[
 \frac vM\ge\frac{4XY^2}{4q^2Y}
       =\frac{wY}{q}\ge wq^2>w.
\]
The resulting $D_*$ equals $Mg$, proving $A_m=0$.

The auxiliary coordinates use the squared-modulus reconstruction already
proved in package~58. For $R=2r_0+1$ with $r_0$ odd, the elementary odd
Pell polynomials satisfy
\begin{gather*}
 \chi_t(2r+1)=tQ_r(t^2),\qquad Q_r(0)=(-1)^r(2r+1),\\
 Q_r(1-A^2)=(-1)^r\psi_A(2r+1).
\end{gather*}
Take $V_* =\chi_I(R)/I$ and $y_* =\psi_I(R)$. They are positive integers,
$V_*>1$, and $I^2V_*^2-(I^2-1)y_*^2=1$. The identities yield
$V_*\equiv-c\pmod{f^2}$ and $V_*\equiv-R\pmod c$.
Since $\gcd(c,f)=1$,
\[
 T=\frac{V_*+c+Rf^2}{cf^2}
\]
is a positive integer and gives the required $V=V_*$. Finally
$(I^2-1)(y_*^2-V_*^2)=V_*^2-1>0$, so $u=y_*-V_*>0$.
All five residuals vanish. This proves both existence implications in
Theorem~\ref{thm:pell26} without changing the fixed tuple.

\subsection{Exact degree and literal cost}

The cancellation identities used in the degree proof are
\begin{align*}
 A_0={}&q^6[r^2+\ell(2r-k)-1]+2q^3(\ell+r)Y+Y^2,\\
 D_{\rm old}={}&X+ac+\gamma H,\qquad L=-v+SX,\\
 A_m={}&M^2[(X+\gamma H)^2+2ac(X+\gamma H)-Hc^2-1]\\
       &\hspace{15mm}+2MD_{\rm old}HL+H^2L^2.
\end{align*}
These are polynomial-ring identities, not relations valid only on solutions.
They give the exact residual degrees
\[
 (\deg A_0,\deg A_m,\deg A_t,\deg A_a,\deg A_s)
       =(11,13,3,13,10).
\]
Put $q_0=3\beta j$, $X_0=x+q_0$, $k_0=\eta+\zeta$, and
$R_4=q_0^3(q_0-F)$. The two highest residual forms are
\begin{align*}
 L_m&=32q_0^2k_0X_0^2Y^4(4q_0^2Y\gamma+R_4),\\
 L_a&=2X_0^2Y^2f^4uR_4.
\end{align*}
Hence the degree-$26$ homogeneous form is $L_m^2+L_a^2$.
The coefficient of $x^4Y^4f^8u^2j^6F^2$ is exactly $4(3\beta)^6>0$;
$L_m^2$ contains no $f$ or $u$ and cannot change it. This establishes
the exact degree for the authentic numerical instance.

The evaluator has $49$ multiplications, $27$ additions and $16$
subtractions, totaling $92$ live binary operations. Scalar multiplication
is charged, every fixed power is expanded into multiplications, and all
fourteen inputs are live. This arithmetic metric is separate from bit
complexity and the height of the expanded coefficients.

\subsection{A degree-30 companion with 85 operations}

Retain the positive inputs
\[
 J,F,\alpha,z,f,i,T,w,Y,r,\eta,\zeta,u,\gamma,
\]
and put $q=\beta J+1$, $X=wq$. Use the same common expressions, but replace
the main carrier and transport by
\begin{align*}
 M&=4qY,\qquad v=k-R-1,\\
 D_*&=M(X+ac)+H(M\gamma-v),\\
 A_t&=(K_0+w)C-F-z(q-1).
\end{align*}
Keep $A_0,A_m,A_a,A_s$ and the five-square sum, calling it $P_{14,30}$.
The literal circuit is in \path{pell/circuits/p14_degree30.json}.

There is no restriction modulo $3$ on $J$ here. Once $q^3\mid Y$, every
prime common to $qY$ and $H$ is $3$. If $3\mid q$, then $27\mid Y$ and
$H\equiv3\pmod{27}$, so its $3$-adic valuation is exactly one. If
$3\nmid q$, then $\gcd(q,H)=1$ and $\gcd(Y,H)\mid3$. As $H$ is odd,
\[
 \gcd(4qY,H)\le3,
 \qquad\frac M{\gcd(M,H)}\ge\frac{4qY}{3}>R+1.
\]
Also $M\mid4XY^2$. The same ordered index proof gives $p=R$ and
$2n=R+1$, then $h=v/E>0$ and $\gamma_{\rm old}=\gamma-v/M>0$.
Completeness uses $\gamma=\gamma_{\rm old}+v/M>0$; both gap
constructions remain unchanged. Thus it has the same consistency
equivalence.

Its residual degrees are $(12,14,2,15,10)$. With
$q_0=\beta J$, $X_0=wq_0$ and $R_4=q_0^3(q_0-F)$, the highest form is
\[
 [P_{14,30}]_{30}=4X_0^4Y^4f^8u^2R_4^2.
\]
The coefficient of $w^4J^{10}Y^4f^8u^2F^2$ is $4\beta^{10}>0$.
The evaluator has $45$ multiplications, $24$ additions and $16$
subtractions, totaling $85$ operations.

\subsection{Verification and remaining formalization}

The independent checker uses the standard Python library and does not
import the circuit builders. It proves three local polynomial identities
by exact sparse algebra, checks all gate dependencies and counts, and
compares the DAGs with separately transcribed formulas at $200$ modular
assignments. It also checks $48$ exact polynomial rays modulo two primes
against the analytic leaders using the actual fixed tuple, exercises
$10{,}584$ bounded carrier cases, and rejects a transcription mutation
for each circuit. The ordinary mathematical degree proof above supplies
the universal upper bound and nonzero leading coefficient; the modular
rays are implementation checks. The same checks pass under Python's
\texttt{-O} mode.

The inherited coefficient audit independently rebuilds the four numerals,
verifies $66{,}062$ finite basis interactions and their provenance, tests
$38{,}663$ one-bit corruptions, and verifies $87{,}380$ small complement
models. These checks do not replace the unbounded native compiler proof.
An independent internal mathematical review of the new carrier is included
as \texttt{pell/INDEPENDENT\_PROOF\_AUDIT.md}; it is AI-assisted review,
not external human peer review.

The next formalization target is the two-stage carrier lemma with its exact
order of recovery: first the modulus $S$, then the index identity $p=R$,
then the full modulus $M$ and $q\mid X$. Potential further reductions must
preserve this order and justify all positive replacement witnesses.
Whether an auxiliary root can remove another coordinate without losing
the degree advantage remains open.


\section{A companion over ten unrestricted integers}
\label{sec:shared-tap-sun}

The shared-kernel-tap source also composes with the unchanged Sun finalizer
of package~59. This gives a fixed polynomial
\[
 F_{10}(f,g,h,k,l,w,x,y,z,m)\in\mathbb Z[f,g,h,k,l,w,x,y,z,m]
\]
of \textbf{exact total degree $7\,428\,120$}, with a division-free evaluator
of \textbf{525 binary arithmetic operations}, such that the inherited
compiler and arithmetic theorems, together with the new source argument,
prove in ZFC
\[
 \neg\exists(f,g,h,k,l,w,x,y,z,m)\in\mathbb Z^{10}\;(F_{10}=0)
 \quad\Longleftrightarrow\quad\operatorname{Con}(\mathrm{ZFC}).
\]
The domain here is all integers. Relative to package~59's ten-integer
polynomial of degree $8\,469\,528$ with 533 gates, the degree decreases by
$12.2959\%$ and eight arithmetic operations are saved. The comparison is
at commit \texttt{beb0192cb8dc36f78c369c080f2adfbc1887a486}.

\subsection{The revised source polynomials}

Use the primary shared-tap certificate, with its complete old31 coefficient
tuple, the original unshifted 31-coordinate positive source, and the two
invertible row operations. The 89 weighted source records use 84 unique
kernel taps. Each shared coefficient is inserted once. The target identities
are
\[
 [Z^{q_j}]\{D_2(Z)U_p(Z)^2+D_1(Z)U_p(Z)+D_0(Z)\}
       =2\rho_j(p),\qquad U_p(Z)=\sum_i p_iZ^{w_i}.
\]
The certified numerical parameters are
\[
 W=116,\quad Q=1841,\quad H=1873,\qquad
 (\deg D_0,\deg D_1,\deg D_2)=(1384,1739,1641).
\]
The leading coefficient of $D_2$ is $+1$. Every kernel coefficient has
absolute value below $C_0=2^{2134927}$. Put
\[
 \kappa=2134940,\qquad \beta_S=2^\kappa>4096C_0.
\]
The even exponent is part of the source binding. Define
\begin{align*}
 b&=1+3f,& B&=\beta_Sb^2,\\
 \ell&=\sum_iB^{w_i},& U_c&=\ell+g,\\
 A_T&=\sum_jB^{q_j},&N_0&=B^{117},\\
 N_1&=B^{1874},&N&=N_0N_1=B^{1991},\\
 M_0&=N_0-1-(b-1)\ell,&J_c&=2G_c+A_T,\\
 S_c&=g+J_cN_0,&T_c&=M_0+(B-2)A_TN_0,\\
 \mathcal C&=bB^{116}.
\end{align*}
Here the remaining expressions are
\begin{align*}
 G_c&=D_2(B)U_c^2+D_1(B)U_c+D_0(B),\\
 X&=N^2+(S_c+T_c)N+N-1-S_c,\\
 Y&=\beta_S^{15}2^{\kappa-21}B^{2114}b
     =2^{2130\kappa-21}b^{4229}.
\end{align*}
These are polynomial abbreviations, with fixed exponents. Their exact
degrees are $7964,4229,233$ for $X,Y,\mathcal C$, respectively.

Suppose first that $b$ is a square, $f\ne0$, and
$0<g<2\mathcal C$. Then $b\ge4$, and
$B^{116}\le U_c<3bB^{116}$. The kernel estimates give
\[
 D_2(B)>\tfrac12B^{1641},\qquad
 |D_1(B)|<2C_0B^{1739},\qquad |D_0(B)|<2C_0B^{1384}.
\]
After normalization by $B^{1641}U_c^2$, the two possibly negative terms
are at most $2C_0/B^{18}$ and $2C_0/B^{489}$, each less than $1/8$.
Consequently $G_c>\tfrac14B^{1641}U_c^2>0$. The upper bound
\[
 |G_c|<26C_0b^2B^{1873},\qquad 2G_c<N_1/64
\]
and the 32-position gap $H-Q$ prove
\[
 0<J_c<N_1/32,\qquad J_c+(B-2)A_T+2<N_1.
\]
Since $0\le M_0<N_0$ and $0<g<N_0$, this yields the strict range
\begin{equation}
 0\le S_c,\qquad 0\le T_c,\qquad S_c+T_c<N.
 \label{eq:new-sun-range}
\end{equation}
It also gives $X\ge N^2>3b$ and $Y>\max(b,256)$. As $\kappa$ is even,
the source congruence is $X\equiv1-22\equiv0\pmod3$.

When $b$ is a power of two and $S_c\mathbin{\&}T_c=0$, the low mask recovers
$g=\sum_i n_iB^{w_i}$ with $0\le n_i<b$. Thus $U_c$ encodes positive
coordinates $p_i=n_i+1$. Every signed convolution coefficient $h_s$ obeys
\[
 |h_s|<993C_0b^2<B/4.
\]
The target padding in $J_c$ and the mask $B-2$ test $h_s=0$ even when
the incoming carry is $-1$. Therefore the same borrow argument proves
\[
 S_c\mathbin{\&}T_c=0
 \quad\Longleftrightarrow\quad
 g\text{ encodes }n_i\in[0,b),\quad \rho_j(n+1)=0\text{ for all }j.
\]
For $b=2^k$ and $B=2^v$, the actual mask population and block exponent are
\[
 s_2(T_c)=139v-31k-22,\qquad \log_2N=1991v.
\]
Hence
\[
 1+\log_2N+s_2(T_c)=2130v-31k-21=\log_2Y.
\]
Using \eqref{eq:new-sun-range}, the exact two-block population identity
and Legendre's formula turn $Y\mid\binom{2X}{X}$ into precisely the source
zero condition. Conversely a positive source zero is not the all-one
vector, since its fifth row has value one there. Choosing arbitrarily
large $b=4^s$ and encoding $n_i=p_i-1$ gives $B\le g<\mathcal C$ and
the required forward source witnesses.

\subsection{The explicit unchanged finalizer}

Here are the fixed polynomial formulas, with local notation distinguished
from the source base and masks:
\begin{gather*}
 L=lY,\quad U=2LX,\quad V=4gwY,\quad W_*=bw,\\
 A=U(V+1),\quad B_*=2X+1,\quad C=B_*+(A-2)h,\\
 P=U^2V,\quad K=X+1+k(P-2),\\
 S_*=2A-5,\quad T_*=3W_*C-2(W_*^2-1),\\
 a_1=(P^2-4)K^2+4,\\
 \Delta_*=8\mathcal C^3gK^2
       -g^2\{32(C-KL)^2\mathcal C^3+g^2K^2\},\\
 O=f^2l^2\Delta_*,\quad \xi=zx,\quad
 a_0=b\{1+(3O-4)b^2z^2\},\\
 D=(A^2-4)C^2+4,\quad E=C^2D\xi,\quad
 F=4(A^2-4)E^2+1,\\
 G_*=2(A+2)(A-2)^2E^2-1,\quad \Omega=16E,\\
 H_*=C+(B_*-C)F+\Omega Fy,\quad
 I_*=(G_*^2-1)H_*^2+1,\quad a_2=DFI_*.
\end{gather*}
Finally set
\begin{align*}
 r&=T_*+S_*y,&\Sigma&=\xi S_*^5,\\
 \Theta&=(r^3-rS_*^2-S_*^3)(r^2+23S_*^2),\\
 w_0&=1,&w_1&=a_0,&w_2&=a_0(P^2+K^2+3),\\
 Z&=\Sigma m-\Theta,&A_i&=\Sigma^2w_i^2a_i\quad(0\le i\le2),\\
 V_n&=Z^2+A_2-A_0-A_1,&
 W_n&=V_n^2+4Z^2A_2-4A_0A_1.
\end{align*}
The promised single polynomial is
\begin{equation}
 \boxed{F_{10}=W_n^2-16V_n^2Z^2A_2.}
 \label{eq:new-shared-sun10}
\end{equation}

The inherited norm, separator, and intersective-quintic argument first
gives $\Sigma\ne0$, all three $a_i$ squares, and $\Sigma\mid\Theta$.
In particular $z\ne0$. These conclusions use no source positivity.
The coprime factors of the square $a_0$ must be signed squares. Since
$b\equiv1\pmod3$, it cannot be a negative square, so $b$ is a positive
square. The case $b=1$ forces $f=O=0$ and $a_0=1-4z^2<0$, impossible.
Thus $b\ge4$. The second factor being a square forces $O>0$, since for
$O\le0$ it is at most $1-4b^2z^2<0$.

Now $O>0$ implies $\Delta_*>0$. As $\mathcal C>0$, its displayed formula
rules out $g\le0$ and $K=0$; division by $gK^2>0$ gives
$g^3<8\mathcal C^3$. Hence $0<g<2\mathcal C$ before the source estimates
are invoked. This establishes the correct order of the signed bootstrap
and explains why no additional positive padding term is needed.

Sun's inherited arithmetic bridge then recovers $b$ a power of two and
the binomial divisibility, giving a positive source zero. The converse
uses the cofinal forward witnesses, the same Pell bridge, and residue
interpolation. The precise dependencies are Sun's 2021 Theorems
4.12--4.13~\cite{sun2021} and their specialization in his 2026
Propositions 6.1--6.2~\cite{sun},
as preserved in the package59 proof. Thus the new part is the verified
source replacement; the compiler and finalizer theorem are inherited.

\subsection{Exact degree certificate and finite checks}

The upper bound is certified by induction at every gate of the literal
division-free circuit: leaves have degrees zero or one, a product uses
the sum of the two bounds, and a sum or difference uses their maximum.
For this circuit the final bound is
\[
 456(7964)+896(4229)+24(233)+1760=7\,428\,120.
\]
An analytic leading-coefficient certificate gives the matching lower
bound. Under the diagonal specialization, the source leading coefficients
are
\begin{align*}
 x_0&=3^{7964}\beta_S^{3982},&
 y_0&=3^{4229}\beta_S^{2129}2^{\kappa-21},&
 c_0&=3^{233}\beta_S^{116}.
\end{align*}
The scaled radicand $A_2$ has degree $1\,857\,030$ and leading coefficient
$2^{362}3^8x_0^{114}y_0^{224}c_0^6$. It uniquely dominates $V_n$, while
$A_2^2$ dominates $W_n$. The final subtracted term has degree
$5\,735\,370<7\,428\,120$, so $A_2^4$ is the leading term of $F_{10}$.
Its diagonal leading coefficient is
\[
 2^{1448}3^{32}x_0^{456}y_0^{896}c_0^{24}
   =2^{3727056\kappa-17368}3^{7426392}>0.
\]
Thus no cancellation can lower the certified degree.

Independently, after setting all ten variables equal to one indeterminate
$t$, the verifier propagates the coefficient of the certified upper degree
through every gate modulo a prime. It uses the actual fixed coefficients,
authenticated against their complete compressed binary binding. The final
leading coefficients are
\[
 \begin{array}{c|ccc}
 p&101&1009&1000003\\\hline
 [t^{7428120}]F_{10}(t,\ldots,t)\bmod p&95&646&515853.
 \end{array}
\]
These residues also agree with the displayed analytic coefficient. Any
one nonzero residue is a second certificate of the matching lower bound.
This is an exact leading-coefficient computation, not a sample equality
of polynomial values. Separately, exact sparse integer arithmetic verifies the source
coefficient identities, all 59 local circuit identities, the Lucas
factorization, and the eight-factor norm identity. Those identity checks
establish that the circuit computes \eqref{eq:new-shared-sun10}.

The machine-readable circuit and replay files are in the companion
\texttt{sun/} component. Its evaluator has 337 multiplications, 161
additions, and 27 subtractions. The verifier imports no circuit builder
and passes both ordinary Python and \texttt{python -O}. Finite borrow,
population, and source-assignment tests supplement the unbounded proofs;
they do not decide consistency or replace the inherited arithmetic
theorems.


\section{Further research questions}
\label{sec:future}
The present results suggest several concrete next problems.
\begin{enumerate}
\item Can a joint search over variable positions, target positions, and
compatible shared coefficients improve on the supplied certificates?
The search must retain all exact source identities and positivity bounds;
height alone is not the objective because final arithmetic cost also
depends on exponentiation chains.
\item Can source equations be transformed by small invertible row
operations to create more equal kernel coefficients without increasing
coefficient heights or breaking the calibrated source interface?
The current construction already shows why private positions are stronger
than correctness requires.
\item Can the two successive congruence recoveries in the integral
carrier be extended to absorb another quotient or another supplied root?
The order is essential: first recover the large modulus for the Pell
index, then use the exact index to recover scale divisibility.
\item Can one retain the fourteen-variable degree bound while reducing
the additional arithmetic gates? Any proposed simplification must address
both sign selection and positivity of reconstructed canonical witnesses.
\item Can the new finite source certificate and integral-carrier proof
be formalized against the active ProveIt interfaces, then joined to the
native compiler and exported as an object-level ZFC proof?
The current closed formal consistency equations are useful reference
endpoints, but cannot be substituted for this task.
\item Can a substantially smaller direct ZFC proof verifier change the
source polynomial enough to beat these constructions without the K215
route? Merely replacing a literal machine by an unspecified universal
polynomial does not provide a comparable explicit equation.
\end{enumerate}

These are research questions, not claimed consequences of the present
certificates. In particular, no eight-natural-variable or nine-integer-
variable result, no universal optimality theorem, and no consistency
decision follows from the improvements reported here.

\section{Reproduction and integration}
\label{sec:reproduction}
The archive contains this paper's complete LaTeX source, its compiled PDF,
the three mathematical components, exact fixed-coefficient recipes and
payloads, source provenance, and an integrated verification entry point.
The top-level \path{README.md} gives the tested commands and identifies
the primary circuits. Run the supplied verification entry point from the
extracted archive; the finite checks do not require executing a ZFC
refutation search. Rebuilding the paper uses two passes of
\texttt{pdflatex -halt-on-error report.tex}.

The repository comparison was frozen before the new work. The complete
snapshot identifier is given in the abstract and in the provenance file.
No remote repository edits are part of this delivery. The components keep
their original source licenses and notices. Newly written material is
provided under the repository-compatible license included in the archive.

\begin{thebibliography}{99}
\raggedright
\bibitem{repo}
V. Reshetnikov and contributors,
\emph{ProveIt: DiophantineConsistency}, commit
\texttt{beb0192cb8dc36f78c369c080f2adfbc1887a486}.
\url{https://github.com/VladimirReshetnikov/ProveIt/tree/beb0192cb8dc36f78c369c080f2adfbc1887a486/Computability/DiophantineConsistency}.
\bibitem{p31}
ProveIt, source 31, \emph{Con(ZFC) in Diophantine form}, source-31
quadratic system and native coefficient binding, at the snapshot in~\cite{repo}.
\bibitem{p44}
ProveIt, source 44, \emph{Sparse B4 encodings}, compact native coefficient
tuple and compiler proof, at the snapshot in~\cite{repo}.
\bibitem{p51}
ProveIt, source 51, \emph{Sparse arithmetic}, squared-modulus auxiliary
quotient and 66-operation construction, at the snapshot in~\cite{repo}.
\bibitem{p54}
ProveIt, source 54, \emph{Calibrated carry encodings}, ten-integer
arithmetic-cost comparison, at the snapshot in~\cite{repo}.
\bibitem{p56}
ProveIt, source 56, \emph{Free Pell coordinates}, free-coordinate and
carry construction, at the snapshot in~\cite{repo}.
\bibitem{p57}
ProveIt, source 57, \emph{Separated degrees and Pell quotients},
at the snapshot in~\cite{repo}.
\bibitem{p58}
ProveIt, source 58, \emph{Integral-root carriers for Con(ZFC)},
at the snapshot in~\cite{repo}.
\bibitem{p59}
ProveIt, source 59, \emph{Degree-separated encodings and an ordinary-root
merger for Con(ZFC)}, at the snapshot in~\cite{repo}.
\bibitem{jones}
J. P. Jones, Universal Diophantine equation,
\emph{Journal of Symbolic Logic} \textbf{47} (1982), 549--571.
\bibitem{bayer}
J. Bayer, M. David, M. Hassler, Yu. Matiyasevich, and D. Schleicher,
\emph{Diophantine equations over $\mathbb Z$: universal bounds and
parallel formalization}, arXiv:2506.20909v2 (2025).
\url{https://arxiv.org/abs/2506.20909v2}.
\bibitem{sun}
Z.-W. Sun, \emph{Ten unknowns for Hilbert's tenth problem over the
integers}, arXiv:2609.01594v2 (2026).
\url{https://arxiv.org/abs/2609.01594v2}.
\bibitem{sun2021}
Z.-W. Sun, Further results on Hilbert's tenth problem,
\emph{Science China Mathematics} \textbf{64} (2021), 281--306.
\url{https://arxiv.org/abs/1704.03504v7}.
\end{thebibliography}
\end{document}
